\subsection{INEQUALITIES BETWEEN INTEGERS}

PROPOSITION 2. Let \(n\) be an integer. Then every cardinal \(\mathfrak{a}\) such that \(\mathfrak{a} \leqslant n\) is an integer. If \(n \neq 0\) , there exists a unique integer \(m\) such that \(n = m + 1\) , and the relation \(a < n\) is equivalent to \(a \leqslant m\) .

If \(\mathfrak{a} \leqslant n\) , there exists a cardinal \(\mathfrak{b}\) such that \(n = \mathfrak{a} + \mathfrak{b}\) (§ 3, no. 6, Proposition 13). Then \((\mathfrak{a} + 1) + \mathfrak{b} = (\mathfrak{a} + \mathfrak{b}) + 1 = n + 1\) (§ 3, no. 3, Proposition 5, Corollary); and since \(n \neq n + 1\) , we have

\[
(a + 1) + b \neq a + b.
\]

Hence \(\mathfrak{a} + 1 \neq \mathfrak{a}\) , which means that \(\mathfrak{a}\) is an integer. If \(n \neq 0\) , we have \(n \geqslant 1\) (§3, no. 2), and therefore there exists a unique cardinal \(m\) such that \(n = m + 1\) (§3, no. 6, Proposition 13 and no. 4, Proposition 8). Since \(m \leqslant n\) , \(m\) is an integer, from what has already been proved. Finally, if an integer \(a\) is such that \(a < n\) , we have \(n = a + b\) , with \(b \neq 0\) (§3, no. 6, Proposition 13); since \(b\) is an integer, we have \(b = c + 1\) and \(n = m + 1 = (a + c) + 1\) . It follows that \(m = a + c\) (§3, no. 4, Proposition 8), hence \(a \leqslant m\) . Conversely, if \(a \leqslant m\) , we have

\[
a \leqslant m + 1 = n;
\]

and if \(a = n = m + 1\) , we would have \(a > m\) , contrary to hypothesis.

COROLLARY 1. Every subset of a finite set is finite.

COROLLARY 2. If X is a subset of a finite set E and X ≠ E, then

\[
\text { Card } (\mathbf {X}) <   \text { Card } (\mathbf {E}).
\]

For X is contained in the complement \(X'\) of a subset of E consisting of a single element; we have \(\text{Card}(X) \leqslant \text{Card}(X')\) and \(\text{Card}(E) = \text{Card}(X') + 1\) , hence (Proposition 2) \(\text{Card}(X') < \text{Card}(E)\) and a fortiori \(\text{Card}(X) < \text{Card}(E)\) .

Definition 1 shows that, conversely, if E is a set such that

\[
\operatorname{Card} (\mathbf {X}) <   \operatorname{Card} (\mathbf {E})
\]

for every subset X of E such that \(X \neq E\) , then E is finite.

COROLLARY 3. If \(f\) is a mapping of a finite set \(\mathbf{E}\) into a set \(\mathbf{F}\) , then \(f(\mathbf{E})\) is a finite subset of \(\mathbf{F}\) .

For Card \((f(\mathbf{E}))\leqslant\) Card (E) (§ 3, no. 2, Proposition 3).

COROLLARY 4. Let E and F be two finite sets with the same number of elements, and let f be a mapping of E into F. Then the following statements are equivalent :

\begin{enumerate}
\def\labelenumi{(\alph{enumi})}
\item
  \(f\) is an injection;
\item
  \(f\) is a surjection;
\item
  \(f\) is a bijection.
\end{enumerate}

It is enough to prove that (a) and (b) are equivalent. If \(f\) is injective, then \(\operatorname{Card}(f(\mathbf{E})) = \operatorname{Card}(\mathbf{E}) = \operatorname{Card}(\mathbf{F})\) , whence \(f(\mathbf{E}) = \mathbf{F}\) (Corollary 2). If \(f\) is not injective, let \(x\) and \(x'\) be two elements of \(\mathbf{E}\) such that \(x \neq x'\) and \(f(x) = f(x')\) . Then, putting \(\mathbf{E}' = \mathbf{E} - \{x\}\) , we have \(f(\mathbf{E}') = f(\mathbf{E})\) , whence \(\operatorname{Card}(f(\mathbf{E})) \leqslant \operatorname{Card}(\mathbf{E}') < \operatorname{Card}(\mathbf{E})\) by virtue of Corollary 2; but since \(\operatorname{Card}(\mathbf{F}) = \operatorname{Card}(\mathbf{E})\) , it follows that \(f(\mathbf{E}) \neq \mathbf{F}\) .

